\documentclass[11pt]{article}
\usepackage[a4paper,margin=1in]{geometry}
\usepackage{amsmath,amssymb,amsthm,mathtools}
\usepackage[hidelinks]{hyperref}
\usepackage{microtype}

\newtheorem{theorem}{Theorem}
\newtheorem{lemma}[theorem]{Lemma}
\newtheorem{remark}[theorem]{Remark}
\DeclareMathOperator{\Residue}{Res}
\DeclareMathOperator{\Log}{Log}

\title{A Boundary Prime--Power Decomposition and\\
Local Perron Residues for Dirichlet $L$-Functions}
\author{Saar Shai}
\date{Proof and novelty audit --- 1 August 2026}

\begin{document}
\maketitle

\begin{abstract}
We record a boundary prime-power decomposition and the exact local residue of
a shifted Perron kernel.  The boundary identity is unconditional and, in fact,
does not use that the evaluation point is a zero: its only delicate term is an
ordinarily convergent prime sum on $\Re z=1$.  We fix the logarithm by a
right-hand Euler-product limit, so no false path-independence assertion is
needed.  The residue is a local Laurent calculation and has been formalized in
Lean~4.  Neither result supplies a global partial-M\"obius-sum asymptotic:
off-target zeros contribute non-decaying oscillatory residues on the Riemann
hypothesis.  A targeted literature audit finds the analytic ingredients in
earlier work of Conrad and Akatsuka; no novelty claim is made for the
decomposition itself.
\end{abstract}

\paragraph{Status.}
This is a proved technical note, not a submission draft.  The theorem and local
residue below are complete at paper level.  The independent novelty gate fails
for a standalone research paper because the theorem is an elementary
specialization of established boundary Euler-product machinery.  No global
$c_K$ asymptotic is claimed.

\section{Setup and branch convention}

Let $\chi$ be a Dirichlet character modulo $q$, let
$s=\tfrac12+i\tau$ with $\tau\ne0$, and let $\psi$ be the primitive
character of conductor $f\mid q$ that induces the (possibly imprimitive)
character $\chi^2$.  Define
\begin{equation}
 T_K(\chi,s)
 :=\sum_{p\le K}\sum_{k\ge2}
       \frac{\chi(p)^k}{k p^{ks}}.
\end{equation}
For $z=1+2i\tau$, define the Euler-product-normalized boundary logarithm
\begin{equation}\label{eq:boundary-log}
 \Log_*L(z,\psi)
 :=\lim_{\epsilon\downarrow0}
   \sum_p\sum_{r\ge1}\frac{\psi(p)^r}{r p^{r(z+\epsilon)}}.
\end{equation}
The limit exists by the boundary Abel theorem used below.  This is a branch
\emph{definition}, not a claim of path independence.  Continuation around a
zero can change a logarithm by $2\pi i$, so arbitrary-path independence would
be false.

Set
\begin{align}
 \mathrm{BPC}_1(\chi,s)
   &:=\frac12\sum_{\substack{p\mid q\\p\nmid f}}
       \Log\!\left(1-\psi(p)p^{-2s}\right),\\
 \mathrm{BPC}_2(\chi,s)
   &:=-\frac12\sum_{r\ge2}\frac1r\sum_p
       \frac{\chi(p)^{2r}}{p^{2rs}},\\
 T_{\ge3}(\chi,s)
   &:=\sum_{k\ge3}\frac1k\sum_p
       \frac{\chi(p)^k}{p^{ks}}.
\end{align}
The last two double sums are absolutely convergent.  In $\mathrm{BPC}_1$,
$|\psi(p)p^{-2s}|=p^{-1}<1$, so $\Log(1-w)=-\sum_{r\ge1}w^r/r$
unambiguously denotes the principal logarithm.

\section{Boundary identity}

\begin{theorem}[Boundary prime-power decomposition]\label{thm:boundary}
Under the setup above, the ordinary prime-cutoff limit exists and
\begin{equation}\label{eq:boundary}
 \lim_{K\to\infty}T_K(\chi,s)
 =\frac12\Log_* L(2s,\psi)
  +\mathrm{BPC}_1(\chi,s)
  +\mathrm{BPC}_2(\chi,s)
  +T_{\ge3}(\chi,s).
\end{equation}
\end{theorem}

\begin{proof}
For every finite $K$, absolute convergence of the inner series gives
\begin{equation}\label{eq:finite-split}
 T_K(\chi,s)=\frac12\sum_{p\le K}\frac{\chi(p)^2}{p^{2s}}
 +\sum_{p\le K}\sum_{k\ge3}\frac{\chi(p)^k}{k p^{ks}}.
\end{equation}
The second term converges absolutely as $K\to\infty$, since
\[
 \sum_p\sum_{k\ge3}\frac{|\chi(p)|^k}{k p^{k/2}}
 \le \frac1{3(1-2^{-1/2})}\sum_p p^{-3/2}<\infty.
\]
Thus no conditional rearrangement is being made.

Put $\vartheta=\chi^2$, viewed as a character modulo $q$.  Conrad's
boundary Dirichlet-series theorem~\cite[Theorem~4.1]{Conrad2005PartialEuler},
together with his comparison of prime and prime-power cutoff orders
\cite[Lemma~3.1]{Conrad2005PartialEuler}, applied to the Euler-product
logarithm of $L(z,\vartheta)$ shows that, for $z=1+2i\tau$ with
$\tau\ne0$, the prime-cutoff series
\begin{equation}\label{eq:full-log-series}
 \sum_p\sum_{r\ge1}\frac{\vartheta(p)^r}{r p^{rz}}
\end{equation}
converges and equals its right-hand boundary value from $\Re z>1$.
The required hypotheses follow from the classical nonvanishing of Dirichlet
$L$-functions on $\Re z=1$; when $\vartheta$ is principal, the only pole is
at $z=1$, excluded by $\tau\ne0$.  Akatsuka independently gives the direct
principal-character estimate
\begin{equation}\label{eq:akatsuka-tail}
 \sum_{X<p\le Y}p^{-1-2i\tau}=O_\tau(1/\log X)
\end{equation}
in equations (2.3)--(2.5) of~\cite{Akatsuka2017EulerProduct}.

The $r\ge2$ part of \eqref{eq:full-log-series} is absolutely convergent.
Subtracting it, without changing order, proves
\begin{equation}\label{eq:prime-isolation}
 \sum_p\frac{\chi(p)^2}{p^{2s}}
 =\Log_*L(2s,\vartheta)
  -\sum_{r\ge2}\frac1r\sum_p
       \frac{\chi(p)^{2r}}{p^{2rs}}.
\end{equation}
Finally, for $\Re z>1$ the inducing-character identity is
\[
 L(z,\vartheta)=L(z,\psi)
 \prod_{\substack{p\mid q\\p\nmid f}}(1-\psi(p)p^{-z}).
\]
Take Euler-product logarithms there and then the horizontal limit to $2s$.
The finite factors stay inside the disk $|w-1|<1$, so their principal logs
agree with their convergent power series.  Hence
\[
 \Log_*L(2s,\vartheta)=\Log_*L(2s,\psi)
 +\sum_{\substack{p\mid q\\p\nmid f}}
   \Log(1-\psi(p)p^{-2s}).
\]
Multiplying \eqref{eq:prime-isolation} by $1/2$ and substituting it in
\eqref{eq:finite-split} proves \eqref{eq:boundary}.
\end{proof}

\begin{remark}[Weakness and convergence mode]
Theorem~\ref{thm:boundary} does not require $s$ to be a zero of
$L(s,\chi)$, nor does it require simplicity.  The limit is ordinary
prime-cutoff convergence.  Only the $k=2$ term is conditional; all other
passages use absolute convergence.  The radial limit in
\eqref{eq:boundary-log} agrees with the ordinary boundary series by Abel's
theorem, but arbitrary complex paths are not part of the statement.
\end{remark}

\section{The local double-pole residue}

\begin{lemma}[Local Perron residue]\label{lem:local-residue}
Let $L$ be holomorphic near $\rho$, with $L(\rho)=0$ and
$L'(\rho)\ne0$.  For $K>1$,
\begin{equation}
 \Residue_{w=0}\frac{K^w}{wL(\rho+w)}
 =\frac{\log K}{L'(\rho)}-
   \frac{L''(\rho)}{2L'(\rho)^2}.
\end{equation}
\end{lemma}

\begin{proof}
Taylor expansion at the simple zero gives
\[
 \frac1{L(\rho+w)}
 =\frac1{L'(\rho)w}-\frac{L''(\rho)}{2L'(\rho)^2}+O(w),
\]
while
\[
 \frac{K^w}{w}=\frac1w+\log K+O(w).
\]
The coefficient of $w^{-1}$ in their product is the claimed residue.
\end{proof}

\begin{remark}[Formal scope]
Lemma~\ref{lem:local-residue} is represented by a zero-\texttt{sorry} Lean~4
theorem for an analytic function with a simple zero at the origin.  The formal
result certifies the local Laurent algebra; it does not certify a Perron
contour shift or Theorem~\ref{thm:boundary}.
\end{remark}

\section{Why no global partial-sum asymptotic is claimed}

Applying a truncated Perron formula to
\[
 c_K(\chi,\rho)=\sum_{n\le K}\frac{\mu(n)\chi(n)}{n^\rho}
\]
and shifting the contour crosses every off-target zero $\rho'$.  In the simple
zero case its contribution contains a factor of the shape
\begin{equation}
 \frac{K^{\rho'-\rho}}
      {(\rho'-\rho)L'(\rho',\chi)}.
\end{equation}
Under RH, $\Re(\rho'-\rho)=0$, so this is oscillatory with modulus independent
of $K$; it is not automatically $o(1)$.  Inoue's explicit formula
\cite{Inoue2021} organizes these zero contributions but does not permit them
to be discarded.  A complete contour displacement would also have to specify
the zero-sum ordering, bound the horizontal and left-edge integrals, and handle
multiple zeros.  None of those global steps is used in
Theorem~\ref{thm:boundary}.  Consequently Lemma~\ref{lem:local-residue}
identifies one local term, not the asymptotic of $c_K$.

\section{Numerical audit design}

A reproducibility computation can use four fixed character/zero pairs and
report, for every cutoff $K$: the left side of \eqref{eq:boundary}; each of
the four right-side components; the complex residual and working precision;
direct checks of $L(\rho)=0$, $L'(\rho)\ne0$, and the inducing character
$\psi$; and executable source, raw values, software versions, and hashes.
Observed rates must be separated from unconditional estimates and from
RH-conditional bounds such as those in Soundararajan~\cite{Soundararajan2009}.

\section{Novelty and disposition}

The exact four-term display in Theorem~\ref{thm:boundary} was not found
verbatim in the targeted search.  That does not make it a new theorem.
Conrad's general boundary framework already identifies the squared-character
second moment and controls the ordering issue
\cite[Example~4.6 and Theorem~4.9]{Conrad2005PartialEuler}; Akatsuka proves the
principal boundary prime-sum estimate directly; and Kaneko treats broader
Dirichlet partial Euler-product asymptotics~\cite{Kaneko2022DirichletEuler}.
The displayed formula is a transparent isolation of the $k=2$ term plus the
standard imprimitive Euler factors.  The local residue is likewise standard
Laurent algebra.  Therefore these two results do not presently clear a
standalone-paper novelty threshold.  Their defensible use is as a technical
appendix and reproducibility certificate supporting a paper with a genuinely
new prime-bias theorem or numerical discovery.

The remaining open work is numerical rather than theorem-closing: rerun the
four character/zero examples from raw inputs, report interval/error control,
and reconcile the branch convention across code and exposition.  Those checks
can strengthen a computational section, but they do not change the novelty
verdict above.

\bibliographystyle{plain}
\bibliography{references}
\end{document}
